\section{硬件辅助 TEE taxonomy}

\begin{frame}[t,shrink=6]{硬件辅助 TEE taxonomy：方向开场}
\DeckKicker{方向开场}
\SlideMeta{证据等级}{方向综述}{来源状态：公开材料综合}
{\large\bfseries\color{Ink} 01 的任务不是罗列 SGX、SEV、TrustZone、CCA、CoVE，而是建立一张能解释 protected object、resource manager、evidence chain 和 device boundary 的总分类图。\par}\vspace{1.8mm}
\begin{columns}[T,totalwidth=\textwidth]
  \begin{column}{0.45\textwidth}
    \fcolorbox{Rule}{Paper}{%
  \begin{minipage}[t]{0.97\linewidth}
    \vspace{1.1mm}
    \hspace{1.4mm}\begin{minipage}{0.92\linewidth}
      \PanelTitle{内容说明}
      \scriptsize \begin{itemize}
  \item TEE 的共同问题是：远程用户不信任平台管理者，但平台仍要调度 CPU、分配内存、处理 I/O，并给 verifier 提供可信证据。
  \item 本方向用 Li 2024 固定通用 server-side TEE runtime taxonomy，用 Boubakri 2025 补 RISC-V 开放 ISA 谱系，用 Wang/Huang 2026 补 accelerator/device TEE 边界。
  \item 三篇都是 SoK/Survey，因此只作为 design-space anchor；具体平台机制、协议状态和性能数字仍回到原始论文或官方规范。
  \item 讲法上先建立分类轴，再进入后续章节的 Arm CCA、RISC-V CoVE/AP-TEE、CoVE-IO、SPDM/TDISP、accelerator TEE。
\end{itemize}
    \end{minipage}
    \vspace{1mm}
  \end{minipage}}
  \end{column}
  \begin{column}{0.49\textwidth}
    \fcolorbox{Rule}{Panel}{%
  \begin{minipage}[t]{0.97\linewidth}
    \vspace{1.1mm}
    \hspace{1.4mm}\begin{minipage}{0.92\linewidth}
      \PanelTitle{01 的分类轴}
      \scriptsize {\tiny
\begin{tabularx}{\linewidth}{@{}YYY@{}}
\toprule
\textbf{分类轴} & \textbf{核心问题} & \textbf{后续章节映射} \\
\midrule
Protected object & 保护 enclave/CVM/Realm/device workload 的哪一段状态 & TrustZone、CCA、CoVE、accelerator TEE \\
Resource manager & host、RTPM/RMM/TSM、TEE instance 或 device controller 谁管资源 & runtime、memory ownership、interrupt、I/O \\
Evidence chain & attestation report 覆盖哪些初始状态和配置 & RATS、CCA token、TVM/device evidence \\
Device boundary & 数据离开 CPU TEE 后是否仍有身份、隔离和加密 & CoVE-IO、SPDM/TDISP、IDE、DPU/GPU \\
\bottomrule
\end{tabularx}
}
    \end{minipage}
    \vspace{1mm}
  \end{minipage}}
  \end{column}
\end{columns}
\vspace{1mm}
{\tiny\textcolor{Muted}{引用依据：本页用三篇 SoK/Survey 建立方向边界，不把它们当作一手机制证据。}}
\end{frame}


\begin{frame}[t,shrink=6]{SoK: Understanding Design Choices and Pitfalls of Trusted Execution Environments：内容摘要}
\DeckKicker{内容摘要}
\SlideMeta{证据等级}{E2 / Peer-reviewed taxonomy anchor}{来源状态：本地 PDF 已核验}
{\large\bfseries\color{Ink} Li 2024 的核心贡献是 TRAF：把 TEE 生命周期和 CPU/memory/I/O runtime event 映射到四种 resource-management mode。\par}\vspace{1.8mm}
\begin{columns}[T,totalwidth=\textwidth]
  \begin{column}{0.45\textwidth}
    \fcolorbox{Rule}{Paper}{%
  \begin{minipage}[t]{0.97\linewidth}
    \vspace{1.1mm}
    \hspace{1.4mm}\begin{minipage}{0.92\linewidth}
      \PanelTitle{内容说明}
      \scriptsize \begin{itemize}
  \item 动机：server-side TEEs 设计差异大，读者只按产品名学习会看不清 host、TCB、TEE instance 的职责边界。
  \item 工作：提出 TEE Runtime Architectural Framework，先分 remote attestation 与 runtime management，再分析 CPU、memory、I/O 事件。
  \item 关键设计：unprotected、RTPM-only、RTPM-guarded、instance-assisted 四种 mode 解释安全、可管理性和 TCB 规模之间的取舍。
  \item 在本报告中，它是 01 的主 taxonomy anchor，但不替代平台论文或官方规范。
\end{itemize}
    \end{minipage}
    \vspace{1mm}
  \end{minipage}}
  \end{column}
  \begin{column}{0.49\textwidth}
    \fcolorbox{Rule}{Panel}{%
  \begin{minipage}[t]{0.97\linewidth}
    \vspace{1.1mm}
    \hspace{1.4mm}\begin{minipage}{0.92\linewidth}
      \PanelTitle{TRAF 分析管线}
      \scriptsize \begin{itemize}
  \item 01. TEE lifecycle: attestation vs runtime
  \item 02. runtime event: CPU / memory / I/O
  \item 03. manager: host / RTPM / TEE instance
  \item 04. mode: unprotected / only / guarded / assisted
  \item 05. implication: security / TCB / manageability
  \item 06. pitfall mapping: design flaw -> known attack
\end{itemize}
    \end{minipage}
    \vspace{1mm}
  \end{minipage}}
  \end{column}
\end{columns}
\vspace{1mm}
{\tiny\textcolor{Muted}{引用依据：机制图按 Li 2024 Figure 2-Figure 5 重绘。}}
\end{frame}


\begin{frame}[t,shrink=6]{SoK: Understanding Design Choices and Pitfalls of Trusted Execution Environments：研究背景}
\DeckKicker{研究背景}
\SlideMeta{证据等级}{E2 / Peer-reviewed taxonomy anchor}{来源状态：本地 PDF 已核验}
{\large\bfseries\color{Ink} TEE 的根本矛盾是：host 被 threat model 视为不可信，却仍必须管理调度、页表、I/O 和实例生命周期。\par}\vspace{1.8mm}
\begin{columns}[T,totalwidth=\textwidth]
  \begin{column}{0.45\textwidth}
    \fcolorbox{Rule}{Paper}{%
  \begin{minipage}[t]{0.97\linewidth}
    \vspace{1.1mm}
    \hspace{1.4mm}\begin{minipage}{0.92\linewidth}
      \PanelTitle{内容说明}
      \scriptsize \begin{itemize}
  \item 论文把 Secure Remote Execution 目标写成 secure measurement、confidentiality、integrity；这些保证依赖 hardware-backed TCB。
  \item 攻击者模型假设 adversary 可控制 privileged software，并管理 non-TEE world 的 memory mapping、I/O devices 和 CPU scheduling。
  \item remote attestation 只能证明初始状态；运行中是否安全，取决于每个 runtime event 的管理权如何被限制。
  \item 商业 TEE 常排除 DoS 和 microarchitectural side-channel，因此页面需要显式标出安全边界。
\end{itemize}
    \end{minipage}
    \vspace{1mm}
  \end{minipage}}
  \end{column}
  \begin{column}{0.49\textwidth}
    \fcolorbox{Rule}{Panel}{%
  \begin{minipage}[t]{0.97\linewidth}
    \vspace{1.1mm}
    \hspace{1.4mm}\begin{minipage}{0.92\linewidth}
      \PanelTitle{三条必须先讲清的边界}
      \scriptsize {\tiny
\begin{tabularx}{\linewidth}{@{}YYY@{}}
\toprule
\textbf{边界} & \textbf{谁参与} & \textbf{为什么影响 taxonomy} \\
\midrule
Trust boundary & host vs TCB vs TEE instance & 决定 attacker 能操作哪些资源 \\
TCB boundary & manufacturer TCB vs ISV TCB & 决定漏洞责任和验证范围 \\
Evidence boundary & attestation/spec/paper & 决定 verifier 可以相信什么 \\
\bottomrule
\end{tabularx}
}
    \end{minipage}
    \vspace{1mm}
  \end{minipage}}
  \end{column}
\end{columns}
\vspace{1mm}
{\tiny\textcolor{Muted}{引用依据：背景页强调 threat model 和 evidence boundary，避免把 SoK 写成平台安全证明。}}
\end{frame}


\begin{frame}[t,shrink=6]{SoK: Understanding Design Choices and Pitfalls of Trusted Execution Environments：解决方案}
\DeckKicker{解决方案}
\SlideMeta{证据等级}{E2 / Peer-reviewed taxonomy anchor}{来源状态：本地 PDF 已核验}
{\large\bfseries\color{Ink} TRAF 的妙处是把“TEE 安不安全”转成“这个 runtime event 由谁管理、谁校验、谁承担 TCB 成本”。\par}\vspace{1.8mm}
\begin{columns}[T,totalwidth=\textwidth]
  \begin{column}{0.45\textwidth}
    \fcolorbox{Rule}{Paper}{%
  \begin{minipage}[t]{0.97\linewidth}
    \vspace{1.1mm}
    \hspace{1.4mm}\begin{minipage}{0.92\linewidth}
      \PanelTitle{内容说明}
      \scriptsize \begin{itemize}
  \item CPU scheduling 多为 unprotected mode，因为云平台需要调度效率；但 host 可通过中断和调度观察 workload。
  \item context switch 多为 RTPM-only mode，因为寄存器和处理器状态必须由 TCB 保存/恢复。
  \item memory allocation、page fault、virtual memory 常在 guarded 与 assisted mode 之间取舍，是 corner case 高发区域。
  \item I/O data transmission 在很多 TEE 中仍依赖 unprotected path 或应用层保护，直接引出 trusted I/O 和 device attestation 章节。
\end{itemize}
    \end{minipage}
    \vspace{1mm}
  \end{minipage}}
  \end{column}
  \begin{column}{0.49\textwidth}
    \fcolorbox{Rule}{Panel}{%
  \begin{minipage}[t]{0.97\linewidth}
    \vspace{1.1mm}
    \hspace{1.4mm}\begin{minipage}{0.92\linewidth}
      \PanelTitle{四种 management mode}
      \scriptsize {\tiny
\begin{tabularx}{\linewidth}{@{}YYYY@{}}
\toprule
\textbf{Mode} & \textbf{谁主导} & \textbf{收益} & \textbf{风险} \\
\midrule
Unprotected & host & 资源管理简单 & 观察和操控面最大 \\
RTPM-only & TCB/RTPM & 安全边界清楚 & TCB 变大、复杂度升高 \\
RTPM-guarded & host + TCB check & 兼顾可管理性 & 校验遗漏形成漏洞 \\
Instance-assisted & TEE instance runtime & 减少 host 观察面 & runtime 复杂度和 DoS 风险 \\
\bottomrule
\end{tabularx}
}
    \end{minipage}
    \vspace{1mm}
  \end{minipage}}
  \end{column}
\end{columns}
\vspace{1mm}
{\tiny\textcolor{Muted}{引用依据：矩阵按 Figure 4 和 Section 4 的 CPU/memory/I/O event 说明重绘。}}
\end{frame}


\begin{frame}[t,shrink=6]{SoK: Understanding Design Choices and Pitfalls of Trusted Execution Environments：实验结果}
\DeckKicker{实验结果}
\SlideMeta{证据等级}{E2 / Peer-reviewed taxonomy anchor}{来源状态：本地 PDF 已核验}
{\large\bfseries\color{Ink} 这是一篇 SoK，无新实验；它能支撑 taxonomy、coverage 和 pitfall framing，不能单独支撑某个平台的安全或性能结论。\par}\vspace{1.8mm}
\begin{columns}[T,totalwidth=\textwidth]
  \begin{column}{0.45\textwidth}
    \fcolorbox{Rule}{Paper}{%
  \begin{minipage}[t]{0.97\linewidth}
    \vspace{1.1mm}
    \hspace{1.4mm}\begin{minipage}{0.92\linewidth}
      \PanelTitle{内容说明}
      \scriptsize \begin{itemize}
  \item 证据对象是 Figure 1-Figure 5、Table 1、Appendix Table 2，以及作者对代表性 TEE 和 known attacks 的整理。
  \item Figure 5 把多种 TEE 的 CPU、memory、I/O event 映射到四种 mode，说明同一分类轴可跨平台比较。
  \item Table 1 把 vulnerable design flaws 与 known attacks 对齐，用于解释 pitfall，不用于证明某平台已修复或已量产。
  \item PPT 中具体机制、协议状态和数值必须回到原始平台论文、官方规范或厂商资料。
\end{itemize}
    \end{minipage}
    \vspace{1mm}
  \end{minipage}}
  \end{column}
  \begin{column}{0.49\textwidth}
    \fcolorbox{Rule}{Panel}{%
  \begin{minipage}[t]{0.97\linewidth}
    \vspace{1.1mm}
    \hspace{1.4mm}\begin{minipage}{0.92\linewidth}
      \PanelTitle{证据使用边界}
      \scriptsize {\tiny
\begin{tabularx}{\linewidth}{@{}YYY@{}}
\toprule
\textbf{证据对象} & \textbf{可支撑} & \textbf{不可支撑} \\
\midrule
Figure 3/4 & runtime event 与 mode 分类 & 单个平台正确性 \\
Figure 5 & 跨平台 design-choice 对比 & 同等级安全排序 \\
Table 1/2 & pitfall 与攻击类型映射 & 修复状态或生产状态 \\
SoK status & 研究框架 & 新系统实验 \\
\bottomrule
\end{tabularx}
}
    \end{minipage}
    \vspace{1mm}
  \end{minipage}}
  \end{column}
\end{columns}
\vspace{1mm}
{\tiny\textcolor{Muted}{引用依据：实验页保守表述：SoK，无新实验。}}
\end{frame}


\begin{frame}[t,shrink=6]{SoK: Understanding Design Choices and Pitfalls of Trusted Execution Environments：文章评价}
\DeckKicker{文章评价}
\SlideMeta{证据等级}{E2 / Peer-reviewed taxonomy anchor}{来源状态：本地 PDF 已核验}
{\large\bfseries\color{Ink} 评价：Li 2024 是 01 最适合的主线语言；作为 SoK，无新实验，落地价值在 checklist，而不是直接给平台背书。\par}\vspace{1.8mm}
\begin{columns}[T,totalwidth=\textwidth]
  \begin{column}{0.45\textwidth}
    \fcolorbox{Rule}{Paper}{%
  \begin{minipage}[t]{0.97\linewidth}
    \vspace{1.1mm}
    \hspace{1.4mm}\begin{minipage}{0.92\linewidth}
      \PanelTitle{内容说明}
      \scriptsize \begin{itemize}
  \item 优势：把产品名转换成 host/RTPM/instance 分工，适合贯穿后续 Arm、RISC-V、I/O 和 accelerator 章节。
  \item 局限：范围偏 server-side general-purpose CPU TEE，不完整覆盖 embedded TrustZone、纯软件 TEE 和 accelerator/device TEE。
  \item 商业化潜力：可转为 cloud TEE 选型 checklist、verifier policy checklist、architecture review 表。
  \item 使用原则：用它解释分类和风险来源；平台机制、部署状态和采购判断仍需一手证据补强。
\end{itemize}
    \end{minipage}
    \vspace{1mm}
  \end{minipage}}
  \end{column}
  \begin{column}{0.49\textwidth}
    \fcolorbox{Rule}{Panel}{%
  \begin{minipage}[t]{0.97\linewidth}
    \vspace{1.1mm}
    \hspace{1.4mm}\begin{minipage}{0.92\linewidth}
      \PanelTitle{文章评价}
      \scriptsize {\tiny
\begin{tabularx}{\linewidth}{@{}YY@{}}
\toprule
\textbf{维度} & \textbf{判断} \\
\midrule
设计优势 & TRAF 抽象清楚，可统一解释 CPU/memory/I/O runtime event \\
主要局限 & 不是平台机制证明，device path 覆盖不足 \\
商业落地 & 高：适合作为架构审查和 verifier policy checklist \\
讲稿定位 & 主 taxonomy anchor，不作性能或生产状态证据 \\
\bottomrule
\end{tabularx}
}
    \end{minipage}
    \vspace{1mm}
  \end{minipage}}
  \end{column}
\end{columns}
\vspace{1mm}
{\tiny\textcolor{Muted}{引用依据：论文原文、官方规范或公开材料；具体结论受证据等级和原文实验范围约束。}}
\end{frame}


\begin{frame}[t,shrink=6]{A Survey of RISC-V Secure Enclaves and Trusted Execution Environments：内容摘要}
\DeckKicker{内容摘要}
\SlideMeta{证据等级}{E2 survey / Background substrate}{来源状态：本地 PDF 已核验}
{\large\bfseries\color{Ink} Boubakri 2025 的作用是把 RISC-V TEE 从 PMP/monitor enclave 串到 CoVE/AP-TEE confidential VM 谱系。\par}\vspace{1.8mm}
\begin{columns}[T,totalwidth=\textwidth]
  \begin{column}{0.45\textwidth}
    \fcolorbox{Rule}{Paper}{%
  \begin{minipage}[t]{0.97\linewidth}
    \vspace{1.1mm}
    \hspace{1.4mm}\begin{minipage}{0.92\linewidth}
      \PanelTitle{内容说明}
      \scriptsize \begin{itemize}
  \item 动机：RISC-V 的开放性带来可定制安全硬件，也带来 secure enclave、TEE、CVM、monitor、runtime 的术语碎片化。
  \item 工作：按 RISC-V primitives、representative systems、comparison matrix 和 open challenges 组织 2016-2025 谱系。
  \item 关键材料：Figure 1 给 timeline，Table 1 横向比较 isolation、memory、secure I/O、TCB、SDK、compliance。
  \item 在 01 中，它是 RISC-V branch 的辅助 SOTA；具体 Keystone/Penglai/SPEAR-V/CoVE/AP-TEE 机制仍回原始材料。
\end{itemize}
    \end{minipage}
    \vspace{1mm}
  \end{minipage}}
  \end{column}
  \begin{column}{0.49\textwidth}
    \fcolorbox{Rule}{Panel}{%
  \begin{minipage}[t]{0.97\linewidth}
    \vspace{1.1mm}
    \hspace{1.4mm}\begin{minipage}{0.92\linewidth}
      \PanelTitle{RISC-V TEE 谱系}
      \scriptsize \textcolor{Accent}{\bfseries 2016 Sanctum}\par\textbf{open-hardware enclave baseline}\par{\tiny\textcolor{Muted}{Figure 1 timeline}}\par\vspace{1.6mm}
\textcolor{Accent}{\bfseries 2020 Keystone}\par\textbf{open enclave framework}\par{\tiny\textcolor{Muted}{Figure 1 timeline}}\par\vspace{1.6mm}
\textcolor{Accent}{\bfseries 2021 Penglai/CURE}\par\textbf{scalable and customizable protection}\par{\tiny\textcolor{Muted}{Table 1 comparison}}\par\vspace{1.6mm}
\textcolor{Accent}{\bfseries 2023 CoVE/AP-TEE}\par\textbf{TVM/TSM confidential VM direction}\par{\tiny\textcolor{Muted}{Section 3.7 and Section 3.14}}\par\vspace{1.6mm}
    \end{minipage}
    \vspace{1mm}
  \end{minipage}}
  \end{column}
\end{columns}
\vspace{1mm}
{\tiny\textcolor{Muted}{引用依据：时间线按 Boubakri 2025 Figure 1 和 Table 1 重绘。}}
\end{frame}


\begin{frame}[t,shrink=6]{A Survey of RISC-V Secure Enclaves and Trusted Execution Environments：研究背景}
\DeckKicker{研究背景}
\SlideMeta{证据等级}{E2 survey / Background substrate}{来源状态：本地 PDF 已核验}
{\large\bfseries\color{Ink} RISC-V TEE 的背景瓶颈是 substrate 丰富但成熟度不齐：PMP/ePMP、privilege stack、MMU、H-extension、MTT/TSM 并不自动组成可量产平台。\par}\vspace{1.8mm}
\begin{columns}[T,totalwidth=\textwidth]
  \begin{column}{0.45\textwidth}
    \fcolorbox{Rule}{Paper}{%
  \begin{minipage}[t]{0.97\linewidth}
    \vspace{1.1mm}
    \hspace{1.4mm}\begin{minipage}{0.92\linewidth}
      \PanelTitle{内容说明}
      \scriptsize \begin{itemize}
  \item 轻量 enclave 通常复用 PMP/ePMP 和 M-mode monitor，部署简单但受 region 数量、granularity 和 monitor 设计限制。
  \item 更强的系统会修改 page table walker、cache/memory metadata 或安全 monitor，带来移植和验证成本。
  \item CoVE/AP-TEE 把问题提升到 confidential VM，需要 TVM、TSM、ABI、memory ownership 和 attestation 组合。
  \item 因此 01 不能把 RISC-V TEE 简化为一个系统，而要先分清 enclave lineage 与 confidential VM lineage。
\end{itemize}
    \end{minipage}
    \vspace{1mm}
  \end{minipage}}
  \end{column}
  \begin{column}{0.49\textwidth}
    \fcolorbox{Rule}{Panel}{%
  \begin{minipage}[t]{0.97\linewidth}
    \vspace{1.1mm}
    \hspace{1.4mm}\begin{minipage}{0.92\linewidth}
      \PanelTitle{RISC-V TEE substrate}
      \scriptsize \begin{itemize}
  \item 01. U/S/M and HS privilege modes
  \item 02. PMP/ePMP memory regions
  \item 03. secure monitor or TSM
  \item 04. virtual memory / hypervisor extension
  \item 05. MTT/GPT/metadata ownership
  \item 06. TVM evidence and verifier decision
\end{itemize}
    \end{minipage}
    \vspace{1mm}
  \end{minipage}}
  \end{column}
\end{columns}
\vspace{1mm}
{\tiny\textcolor{Muted}{引用依据：机制链条按 Section 2.2、CoVE 和 AP-TEE 小节抽象，不替代规范细节。}}
\end{frame}


\begin{frame}[t,shrink=6]{A Survey of RISC-V Secure Enclaves and Trusted Execution Environments：解决方案}
\DeckKicker{解决方案}
\SlideMeta{证据等级}{E2 survey / Background substrate}{来源状态：本地 PDF 已核验}
{\large\bfseries\color{Ink} 这篇 survey 的核心用法是分层：primitive stack 决定隔离基础，system lineage 决定保护对象，open challenges 决定落地缺口。\par}\vspace{1.8mm}
\begin{columns}[T,totalwidth=\textwidth]
  \begin{column}{0.45\textwidth}
    \fcolorbox{Rule}{Paper}{%
  \begin{minipage}[t]{0.97\linewidth}
    \vspace{1.1mm}
    \hspace{1.4mm}\begin{minipage}{0.92\linewidth}
      \PanelTitle{内容说明}
      \scriptsize \begin{itemize}
  \item Primitive stack：privilege mode、PMP/ePMP、MMU、H-extension、crypto 和 memory tagging 是 RISC-V TEE 的底座。
  \item System lineage：Sanctum/Keystone 建立 open enclave，Penglai/CURE/SPEAR-V 展开 memory primitive 分支，CoVE/AP-TEE 走向 confidential VM。
  \item Comparison matrix：Table 1 让读者看到 secure I/O、cache side-channel、implementation validation、SDK、TCB、compliance 的差距。
  \item Open challenges：Section 4 把 fragmentation、programming model、standardization、production validation 作为后续研究入口。
\end{itemize}
    \end{minipage}
    \vspace{1mm}
  \end{minipage}}
  \end{column}
  \begin{column}{0.49\textwidth}
    \fcolorbox{Rule}{Panel}{%
  \begin{minipage}[t]{0.97\linewidth}
    \vspace{1.1mm}
    \hspace{1.4mm}\begin{minipage}{0.92\linewidth}
      \PanelTitle{RISC-V TEE 三层读法}
      \scriptsize {\tiny
\begin{tabularx}{\linewidth}{@{}YYY@{}}
\toprule
\textbf{层次} & \textbf{代表对象} & \textbf{01 中的讲法} \\
\midrule
Primitive & PMP/ePMP、privilege、MMU、H-extension & 决定隔离和 ownership 能力 \\
System lineage & Keystone、Penglai、SPEAR-V、CoVE & 决定保护对象与 TCB 位置 \\
Evidence matrix & Table 1 & 标出 secure I/O、SDK、验证和成熟度缺口 \\
Open challenges & Section 4 & 连接后续标准化和商业落地判断 \\
\bottomrule
\end{tabularx}
}
    \end{minipage}
    \vspace{1mm}
  \end{minipage}}
  \end{column}
\end{columns}
\vspace{1mm}
{\tiny\textcolor{Muted}{引用依据：三层读法把 survey 内容转成 PPT 结构，不引入原文之外的机制结论。}}
\end{frame}


\begin{frame}[t,shrink=6]{A Survey of RISC-V Secure Enclaves and Trusted Execution Environments：实验结果}
\DeckKicker{实验结果}
\SlideMeta{证据等级}{E2 survey / Background substrate}{来源状态：本地 PDF 已核验}
{\large\bfseries\color{Ink} 这是一篇 Survey，无新实验；它能支撑 RISC-V TEE 谱系和差距图，不能支撑单个系统的安全证明或数值结论。\par}\vspace{1.8mm}
\begin{columns}[T,totalwidth=\textwidth]
  \begin{column}{0.45\textwidth}
    \fcolorbox{Rule}{Paper}{%
  \begin{minipage}[t]{0.97\linewidth}
    \vspace{1.1mm}
    \hspace{1.4mm}\begin{minipage}{0.92\linewidth}
      \PanelTitle{内容说明}
      \scriptsize \begin{itemize}
  \item 证据源是 MDPI HTML 与本地 PDF，核心对象是 Figure 1 timeline、Table 1 comparison matrix 和 Section 4 discussion。
  \item Table 1 可直接用于说明 RISC-V 方案在 secure I/O、SDK、TCB、compliance、validation 维度上差异很大。
  \item 论文中转述的系统细节和数值，PPT 只作为阅读地图；若要展示，必须回引 Keystone、Penglai、SPEAR-V、CoVE/AP-TEE 原文。
  \item 这页的输出应是 evidence matrix，而不是把 survey 写成系统评测页。
\end{itemize}
    \end{minipage}
    \vspace{1mm}
  \end{minipage}}
  \end{column}
  \begin{column}{0.49\textwidth}
    \fcolorbox{Rule}{Panel}{%
  \begin{minipage}[t]{0.97\linewidth}
    \vspace{1.1mm}
    \hspace{1.4mm}\begin{minipage}{0.92\linewidth}
      \PanelTitle{证据边界}
      \scriptsize {\tiny
\begin{tabularx}{\linewidth}{@{}YYY@{}}
\toprule
\textbf{可用作} & \textbf{不可用作} & \textbf{处理方式} \\
\midrule
RISC-V lineage map & 系统安全证明 & 回原始论文 \\
Table 1 overview & 性能数值来源 & 回原始图表 \\
standardization gap & 规范最终状态 & 查官方 spec \\
open challenges & 单系统采购判断 & 与厂商/实现证据结合 \\
\bottomrule
\end{tabularx}
}
    \end{minipage}
    \vspace{1mm}
  \end{minipage}}
  \end{column}
\end{columns}
\vspace{1mm}
{\tiny\textcolor{Muted}{引用依据：实验页保守表述：Survey，无新实验。}}
\end{frame}


\begin{frame}[t,shrink=6]{A Survey of RISC-V Secure Enclaves and Trusted Execution Environments：文章评价}
\DeckKicker{文章评价}
\SlideMeta{证据等级}{E2 survey / Background substrate}{来源状态：本地 PDF 已核验}
{\large\bfseries\color{Ink} 评价：Boubakri 2025 是很好的 RISC-V 地图；作为 Survey，无新实验，正确用法是辅助定位，不是证据终点。\par}\vspace{1.8mm}
\begin{columns}[T,totalwidth=\textwidth]
  \begin{column}{0.45\textwidth}
    \fcolorbox{Rule}{Paper}{%
  \begin{minipage}[t]{0.97\linewidth}
    \vspace{1.1mm}
    \hspace{1.4mm}\begin{minipage}{0.92\linewidth}
      \PanelTitle{内容说明}
      \scriptsize \begin{itemize}
  \item 优势：覆盖 2016-2025 RISC-V secure enclave/TEE 谱系，能把 Keystone、Penglai、SPEAR-V、CoVE、AP-TEE 放进同一张图。
  \item 局限：不同系统 threat model、验证平台和成熟度不一致，Table 1 不能被读成同等级安全排名。
  \item 商业化潜力：适合 RISC-V confidential computing 选型入口和 gap analysis。
  \item 落地依赖：AP-TEE/CoVE、IOMMU/IOPMP、SDK、attestation、secure I/O 和厂商生态继续成熟。
\end{itemize}
    \end{minipage}
    \vspace{1mm}
  \end{minipage}}
  \end{column}
  \begin{column}{0.49\textwidth}
    \fcolorbox{Rule}{Panel}{%
  \begin{minipage}[t]{0.97\linewidth}
    \vspace{1.1mm}
    \hspace{1.4mm}\begin{minipage}{0.92\linewidth}
      \PanelTitle{文章评价}
      \scriptsize {\tiny
\begin{tabularx}{\linewidth}{@{}YY@{}}
\toprule
\textbf{维度} & \textbf{判断} \\
\midrule
设计优势 & RISC-V 谱系和比较矩阵清楚 \\
主要局限 & 二手证据，且 draft/spec 状态可能变化 \\
商业落地 & 中：适合 roadmap，依赖标准和生态成熟 \\
讲稿定位 & 辅助 SOTA，连接 08/09/10 方向 \\
\bottomrule
\end{tabularx}
}
    \end{minipage}
    \vspace{1mm}
  \end{minipage}}
  \end{column}
\end{columns}
\vspace{1mm}
{\tiny\textcolor{Muted}{引用依据：评价结合 Section 4、conclusion 与本仓库 E2 survey 边界。}}
\end{frame}


\begin{frame}[t,shrink=6]{SoK: Analysis of Accelerator TEE Designs：内容摘要}
\DeckKicker{内容摘要}
\SlideMeta{证据等级}{E2 / Background substrate}{来源状态：本地 PDF 已核验}
{\large\bfseries\color{Ink} Wang/Huang 2026 的贡献是把 TEE taxonomy 从 CPU/enclave/CVM 推到 GPU、NPU、TPU、FPGA、DPU/SmartNIC 等 accelerator path。\par}\vspace{1.8mm}
\begin{columns}[T,totalwidth=\textwidth]
  \begin{column}{0.45\textwidth}
    \fcolorbox{Rule}{Paper}{%
  \begin{minipage}[t]{0.97\linewidth}
    \vspace{1.1mm}
    \hspace{1.4mm}\begin{minipage}{0.92\linewidth}
      \PanelTitle{内容说明}
      \scriptsize \begin{itemize}
  \item 动机：confidential workload 使用 accelerator 后，数据进入 driver、queue、device memory、I/O bus、firmware/controller 和 scheduler。
  \item 工作：分析 51 个 accelerator TEE studies，提出 Host-type、Acc.-type、Mix-type，并围绕 access control、memory encryption、attestation 组织设计空间。
  \item 关键统计：41/51 studies 在 2022 年之后提出，42/51 依赖 CPU TEE，40/51 支持特定 CPU architecture，43/51 只支持一种 accelerator。
  \item 在 01 中，它补齐 Li 2024 对 device path 覆盖不足的问题。
\end{itemize}
    \end{minipage}
    \vspace{1mm}
  \end{minipage}}
  \end{column}
  \begin{column}{0.49\textwidth}
    \fcolorbox{Rule}{Panel}{%
  \begin{minipage}[t]{0.97\linewidth}
    \vspace{1.1mm}
    \hspace{1.4mm}\begin{minipage}{0.92\linewidth}
      \PanelTitle{Corpus statistics}
      \scriptsize \textbf{studies analyzed}\hfill \textcolor{Accent}{\bfseries 51}\par\textcolor{Accent}{\rule{0.95\linewidth}{1.1mm}}\par{\tiny\textcolor{Muted}{Table I corpus}}\par\vspace{1.4mm}
\textbf{proposed since 2022}\hfill \textcolor{Accent}{\bfseries 41/51}\par\textcolor{Accent}{\rule{0.80\linewidth}{1.1mm}}\par{\tiny\textcolor{Muted}{motivation}}\par\vspace{1.4mm}
\textbf{rely on CPU TEE}\hfill \textcolor{Accent}{\bfseries 42/51}\par\textcolor{Accent}{\rule{0.82\linewidth}{1.1mm}}\par{\tiny\textcolor{Muted}{motivation}}\par\vspace{1.4mm}
\textbf{specific CPU architecture}\hfill \textcolor{Accent}{\bfseries 40/51}\par\textcolor{Accent}{\rule{0.78\linewidth}{1.1mm}}\par{\tiny\textcolor{Muted}{motivation}}\par\vspace{1.4mm}
\textbf{single accelerator type}\hfill \textcolor{Accent}{\bfseries 43/51}\par\textcolor{Accent}{\rule{0.84\linewidth}{1.1mm}}\par{\tiny\textcolor{Muted}{motivation}}\par\vspace{1.4mm}
    \end{minipage}
    \vspace{1mm}
  \end{minipage}}
  \end{column}
\end{columns}
\vspace{1mm}
{\tiny\textcolor{Muted}{引用依据：统计数字来自 SoK corpus，不等同于某个系统的测量结果。}}
\end{frame}


\begin{frame}[t,shrink=6]{SoK: Analysis of Accelerator TEE Designs：研究背景}
\DeckKicker{研究背景}
\SlideMeta{证据等级}{E2 / Background substrate}{来源状态：本地 PDF 已核验}
{\large\bfseries\color{Ink} CPU TEE 只保护 CPU-side execution/memory；一旦 workload offload，敏感数据和控制状态会跨过 device memory、DMA/MMIO、I/O link 和 firmware。\par}\vspace{1.8mm}
\begin{columns}[T,totalwidth=\textwidth]
  \begin{column}{0.45\textwidth}
    \fcolorbox{Rule}{Paper}{%
  \begin{minipage}[t]{0.97\linewidth}
    \vspace{1.1mm}
    \hspace{1.4mm}\begin{minipage}{0.92\linewidth}
      \PanelTitle{内容说明}
      \scriptsize \begin{itemize}
  \item Figure 2 将 accelerator computing 分成 task preparation/termination 和 task computing 两阶段，攻击面贯穿数据准备、传输、计算、释放。
  \item 攻击目标包括 input/output、model parameters、task code、page table、metadata、driver/runtime 和 accelerator environment。
  \item 攻击来源可能是 host OS/hypervisor、malicious task、device/firmware adversary 或 link/physical adversary。
  \item 因此只讲 CPU TEE 的 memory ownership 不够，必须把 device identity、device memory、queue state 和 bus protection 纳入 taxonomy。
\end{itemize}
    \end{minipage}
    \vspace{1mm}
  \end{minipage}}
  \end{column}
  \begin{column}{0.49\textwidth}
    \fcolorbox{Rule}{Panel}{%
  \begin{minipage}[t]{0.97\linewidth}
    \vspace{1.1mm}
    \hspace{1.4mm}\begin{minipage}{0.92\linewidth}
      \PanelTitle{offload path attack surface}
      \scriptsize \begin{itemize}
  \item 01. CVM/enclave task prepares data/model
  \item 02. driver/runtime configures page table and queue
  \item 03. DMA/MMIO crosses I/O bus
  \item 04. accelerator memory stores plaintext or metadata
  \item 05. compute engine executes workload
  \item 06. output returns through host-visible path
\end{itemize}
    \end{minipage}
    \vspace{1mm}
  \end{minipage}}
  \end{column}
\end{columns}
\vspace{1mm}
{\tiny\textcolor{Muted}{引用依据：攻击面流程按 Figure 2 重绘，避免给出可执行攻击步骤。}}
\end{frame}


\begin{frame}[t,shrink=6]{SoK: Analysis of Accelerator TEE Designs：解决方案}
\DeckKicker{解决方案}
\SlideMeta{证据等级}{E2 / Background substrate}{来源状态：本地 PDF 已核验}
{\large\bfseries\color{Ink} Accelerator TEE 至少同时回答三件事：谁能访问设备路径，数据在哪些位置被保护，远程 verifier 如何确认设备身份和配置。\par}\vspace{1.8mm}
\begin{columns}[T,totalwidth=\textwidth]
  \begin{column}{0.45\textwidth}
    \fcolorbox{Rule}{Paper}{%
  \begin{minipage}[t]{0.97\linewidth}
    \vspace{1.1mm}
    \hspace{1.4mm}\begin{minipage}{0.92\linewidth}
      \PanelTitle{内容说明}
      \scriptsize \begin{itemize}
  \item Architecture：Host-type 复用 CPU-side TEE；Acc.-type 把控制逻辑放进 accelerator；Mix-type 组合 CPU-side 与 device-side trust chain。
  \item Access control：CPU TEE、TSM/RMM/firmware、I/O bus filter、accelerator controller 可能共同决定 DMA/MMIO 和 task access。
  \item Memory protection：host memory、device memory、I/O link 和 security metadata 的粒度必须贴合 accelerator access pattern。
  \item Attestation：CPU TEE report 不足以证明设备侧环境，还需要 accelerator HRoT、endorser、reference values 和 verifier policy。
\end{itemize}
    \end{minipage}
    \vspace{1mm}
  \end{minipage}}
  \end{column}
  \begin{column}{0.49\textwidth}
    \fcolorbox{Rule}{Panel}{%
  \begin{minipage}[t]{0.97\linewidth}
    \vspace{1.1mm}
    \hspace{1.4mm}\begin{minipage}{0.92\linewidth}
      \PanelTitle{三类架构 + 三类机制}
      \scriptsize {\tiny
\begin{tabularx}{\linewidth}{@{}YYY@{}}
\toprule
\textbf{维度} & \textbf{关键对象} & \textbf{01 中的作用} \\
\midrule
Architecture & Host / Acc. / Mix & 说明 trust 放在 CPU 还是 device \\
Access control & CVM、firmware、bus、controller & 说明谁能触达 workload 状态 \\
Memory protection & host/device/link/metadata & 说明数据离开 CPU 后的保护 \\
Attestation & HRoT、endorser、reference values & 说明 verifier 如何信任设备 \\
\bottomrule
\end{tabularx}
}
    \end{minipage}
    \vspace{1mm}
  \end{minipage}}
  \end{column}
\end{columns}
\vspace{1mm}
{\tiny\textcolor{Muted}{引用依据：矩阵综合 Figure 1、Table III、Figure 3、Figure 5/6。}}
\end{frame}


\begin{frame}[t,shrink=6]{SoK: Analysis of Accelerator TEE Designs：实验结果}
\DeckKicker{实验结果}
\SlideMeta{证据等级}{E2 / Background substrate}{来源状态：本地 PDF 已核验}
{\large\bfseries\color{Ink} 这是一篇 SoK，无新 end-to-end 系统实验；它的价值是 corpus、taxonomy tables 和设备侧 checklist，而不是替代具体 accelerator TEE 论文或规范。\par}\vspace{1.8mm}
\begin{columns}[T,totalwidth=\textwidth]
  \begin{column}{0.45\textwidth}
    \fcolorbox{Rule}{Paper}{%
  \begin{minipage}[t]{0.97\linewidth}
    \vspace{1.1mm}
    \hspace{1.4mm}\begin{minipage}{0.92\linewidth}
      \PanelTitle{内容说明}
      \scriptsize \begin{itemize}
  \item 证据对象包括 Table I 51-study corpus、Table III security solutions、Table IV/V access-control preference、Table VI-IX memory/attestation analysis。
  \item Table X-XIII 把 TCB size、software stack size、multi-type support 和 plug-and-play support 纳入部署评价。
  \item Figure 4 给出 memory-encryption granularity mismatch 的示例分析，适合说明设计风险，不应扩展成所有设备 TEE 的通用数值结论。
  \item PPT 中对 TDISP、SPDM、PCIe IDE、CoVE-IO 或具体 GPU/DPU 系统的协议/状态判断，必须回到官方规范或原始论文。
\end{itemize}
    \end{minipage}
    \vspace{1mm}
  \end{minipage}}
  \end{column}
  \begin{column}{0.49\textwidth}
    \fcolorbox{Rule}{Panel}{%
  \begin{minipage}[t]{0.97\linewidth}
    \vspace{1.1mm}
    \hspace{1.4mm}\begin{minipage}{0.92\linewidth}
      \PanelTitle{证据边界}
      \scriptsize {\tiny
\begin{tabularx}{\linewidth}{@{}YYY@{}}
\toprule
\textbf{可支撑} & \textbf{不能支撑} & \textbf{处理方式} \\
\midrule
accelerator TEE design space & 单系统安全证明 & 回原始论文 \\
security mechanism taxonomy & TDISP/SPDM/IDE 状态 & 查官方规范 \\
TCB/compatibility lens & 通用性能结论 & 标为 SoK 示例分析 \\
商业 checklist & 厂商支持状态 & 补 vendor/spec 证据 \\
\bottomrule
\end{tabularx}
}
    \end{minipage}
    \vspace{1mm}
  \end{minipage}}
  \end{column}
\end{columns}
\vspace{1mm}
{\tiny\textcolor{Muted}{引用依据：实验页保守表述：SoK，无新实验。}}
\end{frame}


\begin{frame}[t,shrink=6]{SoK: Analysis of Accelerator TEE Designs：文章评价}
\DeckKicker{文章评价}
\SlideMeta{证据等级}{E2 / Background substrate}{来源状态：本地 PDF 已核验}
{\large\bfseries\color{Ink} 评价：这篇 SoK 让 01 看到 device path；作为 SoK，无新实验，商业价值高但必须接上厂商 HRoT、标准接口和云编排证据。\par}\vspace{1.8mm}
\begin{columns}[T,totalwidth=\textwidth]
  \begin{column}{0.45\textwidth}
    \fcolorbox{Rule}{Paper}{%
  \begin{minipage}[t]{0.97\linewidth}
    \vspace{1.1mm}
    \hspace{1.4mm}\begin{minipage}{0.92\linewidth}
      \PanelTitle{内容说明}
      \scriptsize \begin{itemize}
  \item 优势：RQ、corpus、architecture taxonomy 和多张安全机制表格非常适合转成 accelerator TEE checklist。
  \item 局限：51 个 studies 成熟度不一，很多方案只支持单一 accelerator 或特定 CPU，source availability 也有限。
  \item 商业化潜力：confidential AI、GPU cloud、DPU/SmartNIC offload 和 secure device assignment 需求强。
  \item 落地难点：driver/runtime TCB reduction、device attestation、TDISP/SPDM/IDE、IOMMU/SMMU/IOPMP 和 orchestration stack 需要共同成熟。
\end{itemize}
    \end{minipage}
    \vspace{1mm}
  \end{minipage}}
  \end{column}
  \begin{column}{0.49\textwidth}
    \fcolorbox{Rule}{Panel}{%
  \begin{minipage}[t]{0.97\linewidth}
    \vspace{1.1mm}
    \hspace{1.4mm}\begin{minipage}{0.92\linewidth}
      \PanelTitle{文章评价}
      \scriptsize {\tiny
\begin{tabularx}{\linewidth}{@{}YY@{}}
\toprule
\textbf{维度} & \textbf{判断} \\
\midrule
设计优势 & 把设备侧 access/memory/attestation/TCB 放入统一 taxonomy \\
主要局限 & 二手 SoK 证据，不能替代原始系统或规范 \\
商业落地 & 高：confidential accelerator 与 GPU cloud 需求明确 \\
讲稿定位 & 辅助 SOTA，连接 05/10/13/14/15 方向 \\
\bottomrule
\end{tabularx}
}
    \end{minipage}
    \vspace{1mm}
  \end{minipage}}
  \end{column}
\end{columns}
\vspace{1mm}
{\tiny\textcolor{Muted}{引用依据：评价结合 TCB/compatibility tables 与 conclusion，不扩展为具体产品判断。}}
\end{frame}


\begin{frame}[t,shrink=6]{硬件辅助 TEE taxonomy：方向总结}
\DeckKicker{方向总结}
\SlideMeta{证据等级}{方向综述}{来源状态：公开材料综合}
{\large\bfseries\color{Ink} 01 的最终结论：TEE taxonomy 应从 CPU runtime 分工开始，向 RISC-V open-ISA lineage 和 accelerator/device path 扩展，最后用证据等级控制结论强度。\par}\vspace{1.8mm}
\begin{columns}[T,totalwidth=\textwidth]
  \begin{column}{0.45\textwidth}
    \fcolorbox{Rule}{Paper}{%
  \begin{minipage}[t]{0.97\linewidth}
    \vspace{1.1mm}
    \hspace{1.4mm}\begin{minipage}{0.92\linewidth}
      \PanelTitle{内容说明}
      \scriptsize \begin{itemize}
  \item Li 2024 解决总语言：host、RTPM/TCB、TEE instance 如何在 CPU/memory/I/O runtime event 上分工。
  \item Boubakri 2025 补 RISC-V 分支：PMP/monitor enclave、scalable memory protection、CoVE/AP-TEE confidential VM 和标准化缺口。
  \item Wang/Huang 2026 补设备分支：accelerator 的 access control、memory protection、attestation、TCB 和 compatibility。
  \item 仍未解决的问题：SoK 只能固定 design space；后续章节必须用原始论文、官方规范和厂商证据支撑机制、状态和部署判断。
\end{itemize}
    \end{minipage}
    \vspace{1mm}
  \end{minipage}}
  \end{column}
  \begin{column}{0.49\textwidth}
    \fcolorbox{Rule}{Panel}{%
  \begin{minipage}[t]{0.97\linewidth}
    \vspace{1.1mm}
    \hspace{1.4mm}\begin{minipage}{0.92\linewidth}
      \PanelTitle{本方向方法对比}
      \scriptsize {\tiny
\begin{tabularx}{\linewidth}{@{}YYY@{}}
\toprule
\textbf{论文} & \textbf{核心思想} & \textbf{主要局限} \\
\midrule
Li 2024 & TRAF: runtime event -> management mode -> pitfall & 偏 CPU/server-side TEE，不能证明平台机制 \\
Boubakri 2025 & RISC-V primitives -> systems -> open challenges & 二手 survey，状态需跟官方 spec \\
Wang/Huang 2026 & accelerator TEE: access/memory/attestation/TCB & 不能替代 device/system 原始证据 \\
01 takeaway & 分类先行，证据分级约束结论 & 后续章节必须补一手材料 \\
\bottomrule
\end{tabularx}
}
    \end{minipage}
    \vspace{1mm}
  \end{minipage}}
  \end{column}
\end{columns}
\vspace{1mm}
{\tiny\textcolor{Muted}{引用依据：总结页把三篇 SoK/Survey 的互补关系和证据边界放在同一矩阵中。}}
\end{frame}
